% =============================================================================
% CROSS-REFERENCE ANALYSIS: The Alpha Formula and Strong CP Connection
% Date: 2026-03-23
% =============================================================================

\documentclass[11pt]{article}
\usepackage{amsmath,amssymb,amsthm}
\usepackage[margin=1in]{geometry}
\usepackage{tcolorbox}

\newtheorem{theorem}{Theorem}
\newtheorem{proposition}{Proposition}
\newtheorem{corollary}{Corollary}
\theoremstyle{remark}
\newtheorem*{remark}{Remark}

\title{Deep Investigation: How $\bar\theta = 0$ Connects to the $\alpha$ Formula\\
\large The Shared $\mathbb{C}P^2 \times S^3$ Geometry}
\author{DFD Research Programme}
\date{23 March 2026}

\begin{document}
\maketitle

% =============================================================================
\section{Executive Summary}
% =============================================================================

DFD derives both $\alpha^{-1} = 137.036$ and $\bar\theta = 0$ from the \emph{same}
internal manifold $M = \mathbb{C}P^2 \times S^3$.  These are not independent claims
sitting in separate appendices---they are consequences of a single geometric structure.
This document traces the precise connections, identifies the shared topological
ingredients, and analyzes the implications for axion physics.

\begin{tcolorbox}[colback=green!5!white, colframe=green!50!black,
  title=Central Finding]
The fine-structure constant and the strong CP solution share three
topological ingredients:
\begin{enumerate}
\item The K\"ahler structure of $\mathbb{C}P^2$ (controls both $k_{\max} = 60$
  for $\alpha$ and the reality of Yukawa phases for $\bar\theta$)
\item The $S^3$ factor (provides Chern-Simons quantization for $\alpha$
  \emph{and} the $+1$ dimension that makes the mapping torus even for $\bar\theta$)
\item The gauge bundle $E = \mathcal{O}(9) \oplus \mathcal{O}^{\oplus 5}$
  (fixes $k_{\max}$ and encodes the SM hypercharges whose
  $\mathrm{Tr}(Y^2) = 10$ enters both $\alpha$ and the anomaly cancellation
  that protects $\bar\theta$)
\end{enumerate}
\end{tcolorbox}

% =============================================================================
\section{The Strong CP Problem: Standard Physics}
% =============================================================================

The QCD Lagrangian admits a CP-violating topological term:
\begin{equation}
\mathcal{L}_\theta = \bar\theta \,\frac{g_s^2}{32\pi^2}\,
  \mathrm{Tr}(G_{\mu\nu}\tilde G^{\mu\nu}),
\end{equation}
where
\begin{equation}
\bar\theta = \theta_{\rm bare} + \arg\det(M_u M_d).
\end{equation}
The neutron electric dipole moment constrains $|\bar\theta| < 10^{-10}$,
yet the Standard Model provides no reason for this smallness.

The Peccei-Quinn solution introduces an axion field $a(x)$ with decay constant
$f_a \sim 10^{9}\text{--}10^{12}$ GeV, dynamically relaxing $\bar\theta \to 0$
at the cost of a new particle.

% =============================================================================
\section{DFD's Resolution: $\bar\theta = 0$ from Even-Dimensional Mapping Torus}
% =============================================================================

DFD resolves the strong CP problem without introducing any new particles.  The
argument proceeds in two steps (Appendix L of the unified theory):

\subsection{Step 1: Tree-Level CP Invariance}

The microsector $M = \mathbb{C}P^2 \times S^3$ with bundle $E = \mathcal{O}(9)
\oplus \mathcal{O}^{\oplus 5}$ produces:
\begin{itemize}
\item Real Yukawa eigenvalues from the Fubini-Study K\"ahler potential
  $K_{\rm FS} = \log(1 + |z_1|^2 + |z_2|^2)$, which is manifestly real
\item $\arg\det(M_u M_d) < 10^{-19}$ rad (numerically verified)
\item $\theta_{\rm bare} = 0$ because $\mathrm{Tr}(F \wedge F)$
  integrates to a topological integer $8\pi^2 k_3$ on $\mathbb{C}P^2$, not a
  continuous parameter
\end{itemize}
This establishes $\bar\theta = 0$ at tree level.

\subsection{Step 2: All-Orders Protection via Dai-Freed Theorem}

The all-orders upgrade uses the Dai-Freed anomaly formula.  Define the CP
mapping torus:
\begin{equation}
T_{\rm CP} = (M \times [0,1]) / (x,0) \sim (\mathrm{CP}(x), 1).
\end{equation}
The CP anomaly phase is:
\begin{equation}
A_{\rm CP} = \exp\!\left(\frac{i\pi}{2}\,
  \eta(D_{T_{\rm CP}})\right).
\end{equation}
CP is non-anomalous iff $A_{\rm CP} = 1$, i.e.\ $\eta(D_{T_{\rm CP}}) \in 4\mathbb{Z}$.

The critical observation:
\begin{equation}
\dim T_{\rm CP} = \dim M + 1 = 7 + 1 = 8 \quad\text{(even)}.
\end{equation}
On any closed even-dimensional Spin$^c$ manifold, the chirality operator
$\Gamma$ anticommutes with the Dirac operator ($\Gamma D \Gamma^{-1} = -D$),
forcing exact $\pm\lambda$ spectral pairing.  Therefore:
\begin{equation}
\eta(D_{T_{\rm CP}}) = 0, \qquad A_{\rm CP} = 1.
\end{equation}

\begin{theorem}[Strong CP All-Orders Closure]
In the DFD microsector on $M = \mathbb{C}P^2 \times S^3$:
$\bar\theta = 0$ to all loop orders.  No axion is required.
\end{theorem}

% =============================================================================
\section{The $\alpha$ Formula: Same Geometry}
% =============================================================================

The fine-structure constant is derived from Chern-Simons quantization on the
\emph{same} $S^3$ factor, with the level cutoff determined by the \emph{same}
$\mathbb{C}P^2$ and \emph{same} bundle $E$:

\subsection{The Derivation Chain for $\alpha$}

\begin{enumerate}
\item \textbf{Bundle:} $E = \mathcal{O}(9) \oplus \mathcal{O}^{\oplus 5}$
  on $\mathbb{C}P^2$ (hypercharge integrality forces $q_1 = 3$;
  five chiral multiplet types give $n = 5$; minimal-padding gives $a = 9$)

\item \textbf{Bridge Lemma:}
\begin{equation}
k_{\max} = \chi(\mathbb{C}P^2, E) = \chi(\mathcal{O}(9)) + 5\chi(\mathcal{O})
  = 55 + 5 = 60
\end{equation}

\item \textbf{CS quantization on $S^3$:} With SU(2) weight function
$w(k) = \frac{2}{k+2}\sin^2\frac{\pi}{k+2}$ and levels $k = 0, \ldots, 59$:
\begin{equation}
\beta_{U(1)} = \langle k + 2 \rangle = 3.797
\end{equation}

\item \textbf{Hypercharge weighting:}
$w = N_{\rm species}/(g_F \cdot \mathrm{Tr}(Y^2)) = 7/80$

\item \textbf{Microsector trace:} Regular-module branch gives boost
$\varepsilon_{\rm adj} = d^2/(d^2 - 1) = 4096/4095$

\item \textbf{Result:}
\begin{equation}
\boxed{\alpha^{-1} = 137.03599985 \qquad\text{(residual $-0.006$ ppm)}}
\end{equation}
\end{enumerate}

% =============================================================================
\section{The Connection: Shared Geometric Infrastructure}
% =============================================================================

We now identify the precise points where the $\alpha$ derivation and the
$\bar\theta = 0$ proof share geometric ingredients.

\subsection{Ingredient 1: $\mathbb{C}P^2$ K\"ahler Structure}

\begin{center}
\renewcommand{\arraystretch}{1.3}
\begin{tabular}{p{0.45\textwidth}p{0.45\textwidth}}
\hline
\textbf{Role in $\alpha$} & \textbf{Role in $\bar\theta = 0$} \\
\hline
Hirzebruch-Riemann-Roch gives
$\chi(\mathcal{O}(m)) = \binom{m+2}{2}$, which yields
$k_{\max} = 60$ &
Fubini-Study K\"ahler potential is real
$\Rightarrow$ Yukawa eigenvalues are real
$\Rightarrow$ $\arg\det(M_u M_d) < 10^{-19}$ \\[6pt]
$H^4(\mathbb{C}P^2) = \mathbb{Z}$ provides the
topological sector for Spin$^c$ index computation &
$H^4(\mathbb{C}P^2) = \mathbb{Z}$ makes
$\int\mathrm{Tr}(F \wedge F) = 8\pi^2 k_3$ (integer),
so $\theta_{\rm bare}$ is not a continuous parameter \\[6pt]
Canonical Spin$^c$ structure gives $K = \mathcal{O}(-3)$,
hence the $+3$ shift in $d = k + 4$ &
CP involution (complex conjugation in homogeneous
coordinates) is orientation-preserving isometry preserving
Spin$^c$ structure \\
\hline
\end{tabular}
\end{center}

\subsection{Ingredient 2: The $S^3$ Factor}

The $S^3$ factor does \textbf{quintuple duty} (extending the ``quadruple duty''
noted in Appendix L):

\begin{enumerate}
\item \textbf{Generation counting:}
$N_{\rm gen} = 3$ from the index theorem on $\mathbb{C}P^2 \times S^3$
(enters $k_a = 3/(8\alpha)$ through $N_{\rm gen}$)

\item \textbf{Proton stability:}
$\pi_3(S^3) = \mathbb{Z}$ provides conserved baryon number winding

\item \textbf{Gauge emergence:}
$\pi_3(\mathrm{SU}(3)) = \mathbb{Z}$ encoded in $S^3$ topology

\item \textbf{Strong CP solution:}
$\dim S^3 = 3$ contributes the crucial $+1$ to make $\dim T_{\rm CP} = 8$ even

\item \textbf{$\alpha$ derivation:}
Chern-Simons quantization \emph{on $S^3$} with $k_{\max} = 60$ gives
$\alpha^{-1} = 137.036$
\end{enumerate}

Items 4 and 5 establish the direct connection: the \emph{same} $S^3$ that
provides the 3-manifold for Chern-Simons quantization (yielding $\alpha$) also
provides the dimension that forces $\eta = 0$ (yielding $\bar\theta = 0$).

\subsection{Ingredient 3: The Bundle $E = \mathcal{O}(9) \oplus \mathcal{O}^{\oplus 5}$}

\begin{center}
\renewcommand{\arraystretch}{1.3}
\begin{tabular}{p{0.45\textwidth}p{0.45\textwidth}}
\hline
\textbf{Role in $\alpha$} & \textbf{Role in $\bar\theta = 0$} \\
\hline
$\chi(\mathbb{C}P^2, E) = 60 = k_{\max}$ &
$E$ carries the SM fermion content whose CP properties are analyzed \\[6pt]
$\mathrm{Tr}(Y^2) = 10$ enters hypercharge weighting
$w = 7/80$ &
Gravitational-$U(1)_Y$ anomaly cancellation
($\sum d_3 Y = 0$) is required for the Dai-Freed
argument to apply with the SM spectrum \\[6pt]
$a = 9$ from minimal hypercharge twist &
Same hypercharge structure determines which fermion
representations contribute to $\arg\det(M_u M_d)$ \\
\hline
\end{tabular}
\end{center}

\subsection{Ingredient 4: The $(3,2,1)$ Partition}

\begin{center}
\renewcommand{\arraystretch}{1.3}
\begin{tabular}{p{0.45\textwidth}p{0.45\textwidth}}
\hline
\textbf{Role in $\alpha$} & \textbf{Role in $\bar\theta = 0$} \\
\hline
Frame stiffness ratios $\kappa_r = n_r \kappa_0$ give
$k_a = (n_3/n_2) \times 1/(4\alpha) = 3/(8\alpha)$ &
Topological separation: SU(3)$_c$ on $\mathbb{C}^3$
and SU(2)$_L$ on $\mathbb{C}^2$ prevents CKM phases
from feeding into $\theta_{\rm QCD}$ \\[6pt]
Wilson ratio $= (n_2/n_1) \times N_{\rm gen} = 6$
enters lattice verification &
The partition ensures that weak CP violation
(CKM phase $\delta_{\rm CP} \neq 0$) coexists with
strong CP conservation ($\bar\theta = 0$) \\
\hline
\end{tabular}
\end{center}

% =============================================================================
\section{Does DFD Predict $f_a = 0$ (No Axion)?}
% =============================================================================

\begin{tcolorbox}[colback=red!5!white, colframe=red!60!black,
  title=Sharp Prediction: No QCD Axion]
DFD predicts:
\begin{itemize}
\item $\bar\theta = 0$ exactly (to all loop orders)
\item No Peccei-Quinn symmetry is needed or present
\item $f_a$ is undefined---there is no axion field
\item All axion searches (ADMX, ABRACADABRA, CASPEr, IAXO, etc.) will yield null results in the KSVZ/DFSZ parameter window
\end{itemize}
\end{tcolorbox}

The Peccei-Quinn mechanism introduces a global $U(1)_{\rm PQ}$ symmetry whose
spontaneous breaking at scale $f_a$ produces a pseudo-Nambu-Goldstone boson
(the axion).  The axion potential has a minimum at $\bar\theta_{\rm eff} = 0$.

DFD renders this entire mechanism unnecessary: $\bar\theta = 0$ is enforced
by the topology of $\mathbb{C}P^2 \times S^3$, not by a dynamical field.
Therefore:
\begin{itemize}
\item There is no $U(1)_{\rm PQ}$ symmetry in the DFD Lagrangian
\item There is no axion decay constant $f_a$
\item The prediction is not $f_a = 0$ (which would be singular in PQ theory)
  but rather that \emph{the PQ framework does not apply}
\item ``Axion-like particles'' from other sources (string theory moduli, etc.)
  are not excluded, but the QCD axion specifically is absent
\end{itemize}

This is a hard, falsifiable prediction.  Detection of a QCD axion with
coupling $g_{a\gamma\gamma}$ in the standard KSVZ or DFSZ band would
falsify the DFD strong CP mechanism.

% =============================================================================
\section{Can the $\alpha$ Formula Constrain the QCD Vacuum Angle?}
% =============================================================================

This section addresses whether the $\alpha$ derivation chain provides an
\emph{independent} constraint on $\bar\theta$, beyond the mapping-torus
argument.

\subsection{Through $\mathrm{Tr}(Y^2)$}

The hypercharge trace $\mathrm{Tr}(Y^2) = 10$ enters the $\alpha$ formula
through the weighting $w = N_{\rm species}/(g_F \cdot \mathrm{Tr}(Y^2)) = 7/80$.
This same trace is a consequence of the SM fermion content:
\begin{equation}
\mathrm{Tr}(Y^2) = \sum_{\rm SM\,reps} d_3 \cdot d_2 \cdot Y^2 = 10
  \quad\text{(per generation)}.
\end{equation}

The anomaly cancellation condition $\mathrm{Tr}(Y) = 0$ (gravitational-$U(1)_Y$)
is required for the Dai-Freed argument.  Thus:
\begin{quote}
\emph{The same SM spectrum that gives $\mathrm{Tr}(Y^2) = 10$ for $\alpha$ also
satisfies the anomaly cancellation needed for $\bar\theta = 0$.}
\end{quote}
If $\mathrm{Tr}(Y^2)$ were different (exotic fermion content), both $\alpha$
and the strong CP proof could fail simultaneously.

\subsection{Through the Chern-Simons Structure}

The $\alpha$ derivation uses U(1) Chern-Simons theory on $S^3$.
Chern-Simons theory on $S^3$ also governs the QCD vacuum structure:
\begin{itemize}
\item The QCD $\theta$-term is the 4D descendant of the 3D CS invariant:
  $\theta \sim \int_{S^3} \mathrm{Tr}(A \wedge dA + \tfrac{2}{3} A \wedge A \wedge A)$
\item In DFD, the CS level is quantized ($k \in \mathbb{Z}$), and the
  effective coupling is a weighted average over levels $k = 0, \ldots, 59$
\item This quantization means the ``$\theta$-parameter'' in the CS theory is
  not a continuous variable---it is absorbed into the discrete level structure
\end{itemize}

This provides an \emph{independent} perspective on why $\theta_{\rm bare} = 0$:
in the CS formulation on $S^3$, the parameter analogous to $\theta$ is
quantized away.  The CS levels are integers, and the physical coupling
$\alpha$ is determined by the weighted average, leaving no room for a
continuous CP-violating phase.

\subsection{Through the Instanton Number}

On $\mathbb{C}P^2$:
\begin{equation}
\int_{\mathbb{C}P^2} \mathrm{Tr}(F \wedge F) = 8\pi^2 k_3,
  \qquad k_3 \in \mathbb{Z}.
\end{equation}
The same instanton number $k_3 = 1$ that enters the generation counting formula
$N_{\rm gen} = |k_3 \cdot k_2 \cdot q_1| = 3$ also ensures that the topological
charge is integer, not continuous.  Since $k_3$ is fixed by energy minimization
(not by a vacuum angle), $\theta_{\rm bare} = 0$ is automatic.

\subsection{Summary: Three Independent Links}

\begin{center}
\renewcommand{\arraystretch}{1.3}
\begin{tabular}{lll}
\hline
\textbf{Shared Structure} & \textbf{$\alpha$ Role} & \textbf{$\bar\theta$ Role} \\
\hline
$\mathrm{Tr}(Y^2) = 10$ & Hypercharge weighting $w = 7/80$ &
  Anomaly cancellation for Dai-Freed \\
CS on $S^3$ & Quantized levels $\to$ $\beta_{U(1)} = 3.797$ &
  Level quantization absorbs $\theta$ \\
$k_3 = 1$ on $\mathbb{C}P^2$ & Generation count via $N_{\rm gen} = 3$ &
  Instanton charge is integer, not continuous \\
\hline
\end{tabular}
\end{center}

% =============================================================================
\section{Deeper Structural Analysis: Why They Must Be Connected}
% =============================================================================

The connection between $\alpha$ and $\bar\theta = 0$ is not accidental.
It reflects a deeper principle: \textbf{both are outputs of a single
spectral geometry}.

\subsection{The Spectral Action Perspective}

In the noncommutative geometry / spectral action framework that DFD's
microsector inherits from, the gauge kinetic term and the topological term
both arise from the same heat kernel expansion:
\begin{equation}
\mathrm{Tr}\, f(D^2/\Lambda^2) = \sum_n f_n \Lambda^{d-2n} a_{2n}(D^2).
\end{equation}

\begin{itemize}
\item The $a_4$ coefficient contains the Yang-Mills action
  $\mathrm{Tr}(F_{\mu\nu}F^{\mu\nu})$, from which $\alpha$ is extracted
\item The $a_4$ coefficient \emph{also} contains the topological density
  $\mathrm{Tr}(F \wedge F)$, which is the $\theta$-term
\item In DFD's plateau cutoff regularization, the coefficient of the
  topological term is determined by the same geometry that fixes $\alpha$
\end{itemize}

With the plateau cutoff ($f^{(n)}(0) = 0$ for $n \geq 1$), the $a_4$ term
is the dominant contribution.  The gauge kinetic term and the topological term
are \emph{two components of the same $a_4$ coefficient}, evaluated on the
same manifold with the same bundle.

\subsection{The Logical Dependency Graph}

\begin{verbatim}
        CP^2 x S^3  (internal manifold)
            |
    +-------+-------+
    |               |
    v               v
  CP^2 geometry    S^3 topology
    |               |
    +--+--+     +---+---+
    |  |  |     |   |   |
    v  v  v     v   v   v
 k_max Kahler  CS  N_gen dim=3
  =60  real  quant  =3   |
    |  Yukawa  |    |   +1
    |    |     |    |    |
    v    v     v    v    v
  alpha  arg   alpha  k_a  dim(T_CP)
  =1/137 det=0 =1/137 =3/8a  =8 (even)
           |              |       |
           v              v       v
       theta_bare=0    MOND    eta=0
           |          scale      |
           +----------+----------+
                      |
                      v
               theta_bar = 0
              (to all orders)
\end{verbatim}

The diagram shows that $\alpha$ and $\bar\theta = 0$ share the same root
($\mathbb{C}P^2 \times S^3$) and converge through multiple pathways.

% =============================================================================
\section{Contrast with Alternative Solutions}
% =============================================================================

\begin{center}
\renewcommand{\arraystretch}{1.3}
\begin{tabular}{lccc}
\hline
\textbf{Mechanism} & \textbf{New particles?} & \textbf{Predicts $\alpha$?}
  & \textbf{Connection?} \\
\hline
Peccei-Quinn / Axion & Yes ($a$, with $f_a$) & No & None \\
Nelson-Barr & Yes (vector-like quarks) & No & None \\
Massless up quark & No & No & None \\
\textbf{DFD topology} & \textbf{No} & \textbf{Yes} &
  \textbf{Same geometry} \\
\hline
\end{tabular}
\end{center}

DFD is unique in deriving both $\alpha = 1/137$ and $\bar\theta = 0$ from
the same topological structure.  No other proposed solution to the strong CP
problem has any connection to the value of the fine-structure constant.

% =============================================================================
\section{Experimental Discrimination}
% =============================================================================

\subsection{Axion Searches}

\begin{itemize}
\item \textbf{ADMX} (Axion Dark Matter eXperiment): Scanning $\mu$eV range.
  DFD predicts null result.
\item \textbf{ABRACADABRA}: Broadband search. DFD predicts null result.
\item \textbf{CASPEr}: Nuclear magnetic resonance search. DFD predicts null result.
\item \textbf{IAXO} (International Axion Observatory): Solar axion search.
  DFD predicts null result.
\end{itemize}

A detection in the KSVZ/DFSZ coupling band would simultaneously:
\begin{enumerate}
\item Validate the PQ mechanism
\item Falsify DFD's topological strong CP solution
\item Cast doubt on the $\mathbb{C}P^2 \times S^3$ microsector
  (and hence on the $\alpha$ derivation)
\end{enumerate}

Conversely, continued null results from all axion experiments would be
consistent with DFD and would progressively disfavor PQ models.

\subsection{Neutron EDM}

The next-generation nEDM experiments aim for sensitivity
$|d_n| \sim 10^{-28}$ $e\cdot$cm, probing $|\bar\theta| \sim 10^{-13}$.
DFD predicts $d_n = 0$ exactly.  Any nonzero measurement would falsify DFD.

\subsection{The Cross-Check: $\alpha$ Precision}

If DFD's $\alpha$ derivation is correct ($\alpha^{-1} = 137.03599985$ from
the regular-module microsector), then the \emph{same} geometry that produces
this value \emph{must} also produce $\bar\theta = 0$.  This creates a
cross-consistency requirement:
\begin{quote}
\emph{Any modification to the microsector geometry that would allow
$\bar\theta \neq 0$ would simultaneously shift $\alpha$ away from
$1/137.036$.}
\end{quote}
The sub-ppm agreement of $\alpha$ is therefore indirect evidence for the
$\bar\theta = 0$ prediction.

% =============================================================================
\section{Open Questions}
% =============================================================================

\begin{enumerate}
\item \textbf{Quantitative link:} Can one derive a formula of the form
  $|\bar\theta| < f(\alpha)$ that makes the functional dependence explicit?
  The current proof gives $\bar\theta = 0$ exactly, but a perturbative
  bound in terms of $\alpha$ would strengthen the cross-reference.

\item \textbf{Dark matter implications:} If the QCD axion does not exist,
  DFD must provide an alternative dark matter candidate.  The DFD framework
  addresses galactic rotation curves through the $\psi$ field and $\mu$-function
  without particle dark matter, but cluster-scale and CMB constraints require
  further development.

\item \textbf{Axion-like particles:} DFD excludes the \emph{QCD} axion
  (the one that solves strong CP) but does not necessarily exclude all
  axion-like particles from other sources.  The distinction should be
  made precise.

\item \textbf{Higher-order Seeley-DeWitt coefficients:} The $a_4$ coefficient
  contains both the gauge kinetic and topological terms.  Does the plateau
  cutoff (which eliminates $a_6, a_8, \ldots$ contributions to $\alpha$) also
  provide additional protection for $\bar\theta$ beyond the mapping-torus
  argument?

\item \textbf{Lattice verification:} The lattice Monte Carlo that confirms
  $\alpha = 1/137$ at $k_{\max} = 60$ could in principle be extended to
  measure the topological susceptibility $\chi_t$ on the discretized
  $\mathbb{C}P^2 \times S^3$.  A vanishing $\chi_t$ would provide numerical
  confirmation of $\bar\theta = 0$.
\end{enumerate}

% =============================================================================
\section{Conclusions}
% =============================================================================

\begin{tcolorbox}[colback=blue!5!white, colframe=blue!50!black,
  title=Summary of Findings]

\textbf{1. The $\alpha$ formula and $\bar\theta = 0$ are geometrically entangled.}
Both emerge from the single internal manifold $\mathbb{C}P^2 \times S^3$ with
bundle $E = \mathcal{O}(9) \oplus \mathcal{O}^{\oplus 5}$.

\textbf{2. The $S^3$ factor plays a dual role:}
\begin{itemize}
\item For $\alpha$: provides the compact 3-manifold for CS quantization
\item For $\bar\theta$: contributes $+1$ to $\dim T_{\rm CP}$, making it even (8)
\end{itemize}

\textbf{3. DFD predicts no QCD axion.}
The PQ mechanism is unnecessary; $f_a$ is undefined.  All KSVZ/DFSZ axion
searches should yield null results.

\textbf{4. Three independent cross-links:}
$\mathrm{Tr}(Y^2)$, CS level quantization, and instanton number $k_3$
each connect the $\alpha$ derivation to the $\bar\theta$ proof.

\textbf{5. Mutual falsifiability:}
Any experimental result that falsifies $\bar\theta = 0$ (axion detection,
nonzero nEDM) would simultaneously undermine the $\alpha$ derivation, since
both rely on the same geometry.  Conversely, the sub-ppm success of $\alpha$
is indirect evidence for $\bar\theta = 0$.
\end{tcolorbox}

\end{document}
